\documentclass[11pt]{article}
\usepackage[letterpaper,margin=0.68in,headheight=15pt]{geometry}
\usepackage[T1]{fontenc}
\usepackage{lmodern}
\usepackage{amsmath,amssymb}
\usepackage{array,booktabs,enumitem,multicol}
\usepackage{xcolor}
\usepackage{tikz}
\usepackage{fancyhdr}
\usepackage{microtype}
\setlength{\parindent}{0pt}
\setlength{\parskip}{5pt}
\definecolor{olive}{HTML}{555B35}
\definecolor{gold}{HTML}{B18A22}
\definecolor{soft}{HTML}{F4F1DE}
\definecolor{bluegray}{HTML}{E8EEF2}
\pagestyle{fancy}
\fancyhf{}
\fancyhead[L]{\textcolor{olive}{\small MATH\& 141: Precalculus I}}
\fancyhead[R]{\textcolor{gold}{\small Fall 2026}}
\fancyfoot[C]{\thepage}
\renewcommand{\headrulewidth}{0.45pt}
\renewcommand{\footrulewidth}{0pt}
\setlist[enumerate]{leftmargin=*,itemsep=0.08in,topsep=0.05in}
\newcommand{\assessmenttitle}[2]{%
  {\Large\bfseries Module #1: #2}\par\vspace{3pt}
}
\newcommand{\studentline}{%
\noindent\textbf{Name:}\hspace{2.6in}\textbf{Date:}\hspace{1.15in}\par\vspace{7pt}}
\newcommand{\directions}[1]{\textit{#1}\par\vspace{4pt}}
\newcommand{\answerblank}[1]{\vspace{#1}}
\newcommand{\blankgrid}[4]{%
\begin{center}
\begin{tikzpicture}[x=0.75cm,y=0.75cm]
  \draw[step=1,gray!28,very thin] (#1,#3) grid (#2,#4);
  \draw[->,thick] (#1,0)--(#2,0) node[right] {$x$};
  \draw[->,thick] (0,#3)--(0,#4) node[above] {$y$};
  \foreach \x in {#1,...,#2}{\ifnum\x=0\else\draw (\x,0.10)--(\x,-0.10) node[below=2pt,font=\scriptsize]{\x};\fi}
  \foreach \y in {#3,...,#4}{\ifnum\y=0\else\draw (0.10,\y)--(-0.10,\y) node[left=2pt,font=\scriptsize]{\y};\fi}
\end{tikzpicture}
\end{center}}
\newcommand{\smallgrid}[4]{%
\begin{tikzpicture}[x=0.50cm,y=0.50cm]
  \draw[step=1,gray!28,very thin] (#1,#3) grid (#2,#4);
  \draw[->,thick] (#1,0)--(#2,0) node[right] {$x$};
  \draw[->,thick] (0,#3)--(0,#4) node[above] {$y$};
\end{tikzpicture}}
\begin{document}

% MODULE 2 PREQUIZ
\assessmenttitle{2}{Lines and Parabolas -- Prequiz}
\studentline
\directions{Four short questions. Show the algebra that supports each answer.}
\begin{enumerate}
\item Find the slope of the line through $(-1,2)$ and $(3,10)$.
\answerblank{0.45in}
\item Find an equation of the line with slope $-3$ that passes through $(2,5)$.
\answerblank{0.55in}
\item For
\[
y=2(x+4)^2-1,
\]
state the vertex and tell whether the parabola opens upward or downward.
\answerblank{0.48in}
\item Complete the square to rewrite
\[
x^2+6x+1
\]
in vertex form.
\answerblank{0.72in}
\end{enumerate}
\newpage

% MODULE 2 CHECKIN PAGE 1
\assessmenttitle{2}{Lines and Parabolas -- Weekly Check-In}
\studentline
\directions{Three questions. Show the algebra that supports each answer.}
\textbf{1. A line from two points.}

A line passes through $(-2,1)$ and $(4,10)$.
\begin{enumerate}[label=(\alph*),leftmargin=0.35in]
\item Find the slope.
\answerblank{0.42in}
\item Find an equation of the line in slope-intercept form.
\answerblank{0.65in}
\end{enumerate}

\textbf{2. Completing the square.}

Consider
\[
q(x)=x^2-6x+5.
\]
Complete the square to rewrite $q(x)$ in vertex form. Then state the vertex and whether the parabola opens upward or downward.
\answerblank{1.45in}
\newpage

% MODULE 2 CHECKIN PAGE 2
\assessmenttitle{2}{Lines and Parabolas -- Weekly Check-In, continued}
\studentline
\textbf{3. Transforming and graphing a parabola.}

Begin with the parent function $f(x)=x^2$ and define
\[
g(x)=-2f(x-1)+3.
\]
\begin{enumerate}[label=(\alph*),leftmargin=0.35in]
\item Describe the transformations from $f$ to $g$ in words.
\answerblank{0.48in}
\item Write $g(x)$ without function notation and state its vertex.
\answerblank{0.48in}
\item Graph $g(x)$ below. Clearly mark the vertex and at least two symmetric pairs of points.
\end{enumerate}
\blankgrid{-3}{6}{-6}{6}

\end{document}
